\section{Introduction}
\label{sec:introduction}

\subsection{The Stationary Distribution Problem in Standard ReCom}

Redistricting ensemble methods rely on Markov chain Monte Carlo to generate
a diverse collection of valid redistricting plans.
The seminal ReCom proposal~\citep{deford2021} merges two adjacent districts,
samples a uniformly random spanning tree of the merged region, and accepts the
first balanced cut found in that tree.
The procedure is fast---$O(m \log m)$ per step where $m$ is the number of
edges in the merged region---and produces plans with high acceptance rates.

However, standard \recom{} has a fundamental distributional gap:
the acceptance ratio is implicitly~1 at every step.
Because the forward and reverse proposals are symmetric under the uniform
spanning tree measure, cancelling the ratio exactly, the Metropolis-Hastings
condition $Q(\pi, \pi') / Q(\pi', \pi) = \pi^*(\pi') / \pi^*(\pi)$ cannot
distinguish any two plans.
The chain therefore targets some stationary distribution, but that
distribution is not the uniform distribution over valid plans---it is an
unknown distribution shaped by the spanning tree measure and the geometry
of the adjacency graph.
For litigation contexts where ensemble arguments are used to demonstrate that
a submitted plan is an outlier, this unknown distribution is a significant
weakness.

\subsection{Two Solutions at Different Costs}

Two approaches restore explicit distributional targeting.

\paragraph{Forest ReCom~\citep{mccartan2020,dellaLibera2026forestRecom}.}
Sample a spanning \emph{forest} (rather than a spanning tree) of the merged
region; the acceptance ratio is the ratio of the number of spanning forests
of the two proposed sub-plans, which equals the ratio of matrix-tree
determinants.
This exactly targets the uniform distribution over valid spanning-forest
partitions.
The cost is $O(m^2)$ per step due to the determinant computation.

\paragraph{Merge-Split MCMC~\citep{janson2022} (this paper).}
Sample an independent reverse spanning tree of the same merged region and
count the number of balanced cuts in each tree.
The acceptance ratio is $\mathrm{cuts\_fwd} / \mathrm{cuts\_rev}$, which is
a random estimate of the exact Metropolis-Hastings ratio.
The cost is $O(m \log m)$ per step---the same as standard \recom{}---because
both trees are sampled via Wilson's algorithm~\citep{wilson1996}.

Merge-Split approximates the uniform distribution in expectation: the ratio
is an unbiased estimate of the true MH ratio, so the chain is ``approximately
reversible'' rather than exactly reversible.
The approximation is the price of the $O(m \log m)$ cost; Forest \recom{}
pays $O(m^2)$ to remove the approximation.

\subsection{Main Contributions}

\begin{enumerate}
  \item \textbf{Algorithm.} We give the complete specification of
        \mergesplit{} as implemented in \texttt{redist-cli}
        (\texttt{MergeSplitChain}), including the deterministic dual-seed
        schedule and the two-tree acceptance ratio
        (Section~\ref{sec:algorithm}).

  \item \textbf{Formal properties.} We state the approximate detailed balance
        theorem for \mergesplit{} (Theorem~\ref{thm:approx-balance}), prove
        that the acceptance rate is positive on well-connected graphs
        (Proposition~\ref{prop:positive-accept}), and derive the per-step
        cost (Proposition~\ref{prop:cost})
        (Section~\ref{sec:algorithm}).

  \item \textbf{Ratio direction.} We derive the Metropolis-Hastings ratio
        $Q(\pi' \!\to\! \pi) / Q(\pi \!\to\! \pi')$ carefully, showing that
        the forward cut count appears in the numerator (not the denominator)
        and verifying that the implementation uses
        $\mathrm{cuts\_fwd} / \mathrm{cuts\_rev}$ correctly
        (Section~\ref{sec:algorithm}).

  \item \textbf{Comparison.} We provide a unified three-way comparison of
        standard \recom{}, \mergesplit{}, and \forestrecom{} across
        stationary distribution, acceptance ratio form, per-step cost, and
        empirical acceptance rate (Section~\ref{sec:comparison}).

  \item \textbf{Usage guidance.} We specify when to prefer \mergesplit{}
        over \forestrecom{} and SMC, and provide the CLI invocation
        (Section~\ref{sec:when-to-use}).
\end{enumerate}

\subsection{Paper Structure}

Section~\ref{sec:background} reviews standard \recom{}, explains why its
acceptance ratio is always~1, and introduces the independence argument that
Merge-Split uses to break the symmetry.
Section~\ref{sec:algorithm} defines \mergesplit{} formally, proves its
properties, and derives the ratio direction.
Section~\ref{sec:comparison} compares the three main reversible-ReCom designs.
Section~\ref{sec:when-to-use} gives practical usage guidance.
Section~\ref{sec:conclusion} concludes.
